\section{Evaluation Design}

\subsection{Open-loop controlled benchmark}

The benchmark contains a damped oscillator, forced Duffing oscillator,
controlled Van der Pol oscillator, torque-driven pendulum, and cart--pole.
Fourth-order Runge--Kutta integration generates training, validation, and test
trajectories. Each of eight seeds independently draws the data and initializes
training. Methods share each seed--system data set.

Three test regimes are fixed: clean in-distribution (ID), corrupted ID, and
corrupted out-of-distribution (OOD) parameters. Corruption combines additive
measurement noise and channel missingness with causal carry-forward. Eleven
variants include recurrent, transformer, neural ODE, physics-guided, nominal,
and residual predictors plus filter-gain and training-horizon sensitivities.
Every learned method receives the same six-trial validation budget. The primary
metric is normalized root mean squared error at the 25-step horizon
(NRMSE-25), averaged equally across the 15 system--regime cells within each
seed.

The predeclared primary contrast compares \method{} with the PGNN-style model.
It requires a paired 95\% interval lower bound above a 5\% improvement and an
exact one-sided seed sign-flip $p<0.05$. The gated robustness claim additionally
requires improvement in the two corrupted regimes and noninferiority within
5\% on clean ID data. Component analyses compare the fixed filter plus nominal
rollout with the complete model, vary the residual multiplier at fixed weights,
and impose a synchronous two-step context outage.

\subsection{Frozen closed-loop regulation study}

The additional control study was specified before execution and reuses the
frozen PGNN-style and \method{} checkpoints. It includes all five systems,
eight seeds, and three plant/observation conditions: ID clean, ID corrupt, and
OOD-parameter corrupt. Four independently generated plant episodes populate
each seed--system--condition cell. Initial states, plant parameters,
observation draws, masks, and planner draws are paired across four prediction
interfaces: nominal propagation from the last imputed state, filter plus
nominal propagation, PGNN-style, and \method. ID clean and ID corrupt also
share the plant draw.

Seven zero-input transitions establish the eight-state context. The controller
then regulates for 20 steps. Deterministic-mean cross-entropy MPC uses horizon
8, 64 candidates, two iterations, and eight elites. Action sequences start
with standard deviation 0.6 of the action limit, with a 0.05 floor. The stage
objective is mean squared state error in training-standard-deviation units plus
0.01 times squared normalized action; terminal state error receives an
additional factor of two. The pendulum target is
$(\sin\theta,\cos\theta,\omega)=(0,1,0)$ and the other targets are zero.
Corrupt observations use Gaussian noise with scalar standard deviation 8\% of
the mean training-state standard deviation and 25\% independent channel
missingness.

The primary closed-loop response is the realized mean normalized state loss
over 20 controlled transitions. Episodes, systems, and conditions receive
equal weight within each seed. The reported \method{}--PGNN contrast is the
mean of eight seed-level relative improvements, with a 5,000-replicate paired
seed bootstrap interval and an exact one-sided sign-flip test against zero.
Controller settings were not changed after observing outcomes.

\subsection{Measured-flight external evaluation}

The NanoDrone experiment uses recorded states and motor commands from the
nonlinear system-identification benchmark \cite{busetto2026_nanodrone}. Three
held-out flights from a distinct trajectory family provide 387 forecast
windows. A reduced quadrotor prior omits parts of the recorded platform
dynamics; learned models have 47,000 parameters and use five training seeds.
The evaluation crosses four context conditions (recorded, added noise,
missingness, and both), two horizons (25 and 50), and four predictors
(\method, its unfiltered residual counterpart, a recurrent model, and the
nominal prior). It reports full-state NRMSE, position RMSE, velocity RMSE, and
marginal 95\% coverage. Flights receive equal weight before averaging training
seeds.

Training, test, and analysis artifacts are hash-checked before tables and
figures are generated. The supplementary material gives the system equations,
parameter ranges, tuning space, complete cell results, and artifact inventory.

